% ================================================================
\section{Project Overview \& Goals}
\label{sec:overview}
% ================================================================

The Python package \texttt{fleet\_management} schedules a fleet of $F$ trains over a
horizon of $T$ days. Every day, $M$ missions must be flown. Each train flies at most
one mission, spends the day in the depot, or stays idle. Each train has $L$
components, and each component degrades stochastically. On a depot day a component
may be repaired (imperfect repair) or replaced. The schedule is a mixed-integer
program solved with Gurobi.

\textbf{Goal.} Minimise the cost of depot days, repairs and replacements, plus a
penalty on the largest train damage (Section~\ref{sec:objective}), under a
reliability requirement: for every component and every day,
$\Pr(D_{i\ell k} > \tau_{i\ell}) \leq \varepsilon_{i\ell}$, with $D_{i\ell k}$ the
degradation and $\tau_{i\ell}$ the failure threshold. The horizon has two phases. The
transitory phase ($H_1$ days) starts from the current state of the fleet. The
operating phase ($H_2$ days) is the schedule to repeat. The requirement holds on every
day of the horizon and of every later repetition of the operating phase
(Section~\ref{sec:loop}).

\textbf{Degradation models.} Each cell (component $\ell$ of train $i$) has its own
model, chosen by the input \texttt{model}:
\begin{itemize}[nosep,leftmargin=1.5em]
  \item \textbf{Gaussian}: a Gaussian random walk; tracks the mean and the variance,
        and a repair reduces both (Section~\ref{sec:dyn_gaussian}).
  \item \textbf{Inverse Gaussian (IG)}: tracks the mean; reliability by a cap on the
        mean (Section~\ref{sec:dyn_ig}).
  \item \textbf{Wiener}: Brownian motion with drift; tracks the mean and the
        variance. The variance grows by $\sigma^2$ per mission day; a repair acts
        on the mean only, and only a replacement lowers the variance
        (Section~\ref{sec:dyn_wiener}).
  \item \textbf{Gamma}: Gamma increments with mission-dependent rates; tracks the
        mean; reliability by a cap that is safe for every mix of rates
        (Section~\ref{sec:dyn_gamma}).
  \item \textbf{Rainflow}: distribution-free. The increment law is measured by
        rainflow counting of load histories and enters only through summaries
        (mean, variance, support width, cumulant generating function). Reliability
        is certified by one of five tail bounds: Markov, Cantelli, Hoeffding,
        Bernstein or Chernoff (Section~\ref{sec:dyn_rainflow}).
\end{itemize}
Both the Gaussian and the Wiener model use Gaussian increments. The Gaussian model
tracks a variance that a repair reduces. The Wiener model fits components whose
uncertainty does not depend on the route and is not reduced by a repair. When in
doubt, use the Gaussian model. The rainflow model makes no distributional
assumption: its certificate holds for every increment law with the measured
summaries, at the price of a looser bound than an exact distribution function.

\textbf{Repair and replacement.} Each cell has a repair law, $\ARDinf$ (the default)
or $\ARDone$ (Section~\ref{sec:repair_laws}); every model accepts both. Each cell may
be replaced (Section~\ref{sec:replacement}).
The loop row forces at least one replacement in the operating phase in two cases.
This is not an input check; a schedule without the replacement is infeasible.
\begin{itemize}[nosep,leftmargin=1.5em]
  \item A Wiener cell whose train flies at least one mission in the operating phase:
        its variance falls only by replacement.
  \item An $\ARDone$ cell with $\rho_{i\ell} < 1$ that accumulates mean damage in the
        operating phase: a repair never removes the share $1 - \rho$ of the damage
        since the last intervention (Proposition~\ref{prop:a_ard1_replacement}).
\end{itemize}

\textbf{Problem class.} The problem is a mixed-integer linear program (MILP) when every
reliability row is linear and the damage penalty is \texttt{inf\_norm}. Otherwise it
is a mixed-integer quadratically constrained program (MIQCP). It is nonconvex when a
quadratic reliability row uses the \texttt{exact} form
(Section~\ref{sec:linear_forms}); it is convex when the only quadratic row is the
\texttt{quadratic} damage penalty (Section~\ref{sec:damage_reg}).

\textbf{Scope of the package.} The package solves one schedule (\texttt{solve}),
sweeps the horizon lengths (Section~\ref{sec:sweeps}), starts a solve from a
previous solution or a heuristic (Section~\ref{sec:warm_starts}), plots a schedule
(Section~\ref{sec:plotting}), and checks a result against its input
(Section~\ref{sec:validation}). It reads inputs in YAML, JSON or HDF5
(Section~\ref{sec:schema}) and writes the outputs of Section~\ref{sec:output}.

\textbf{Authority.} This specification is authoritative for the formulation and for
the tools. The code follows it: every change is written here first.

\textbf{Non-goals.}
\begin{itemize}[nosep,leftmargin=1.5em]
  \item Routing decisions are not optimised: the missions are inputs.
  \item Continuous-time formulations are out of scope: all dynamics are
        discrete-time, with a step of one day.
\end{itemize}

% ================================================================
\section{Index Sets and Notation}
\label{sec:notation}
% ================================================================

\begingroup\small
\begin{longtable}{@{}p{4.6cm}p{8.2cm}p{2.2cm}@{}}
\toprule
Symbol & Meaning & Defined in \\
\midrule
\endhead
\multicolumn{3}{@{}l}{\emph{Indices and horizon}} \\
$i \in \{1, \ldots, F\}$ & train & \\
$\ell \in \{1, \ldots, L\}$ & component of a train & \\
$(i, \ell)$ & a \emph{cell}: component $\ell$ of train $i$ & \\
$j \in \{0, 1, \ldots, M\}$ & activity: $j = 0$ is a depot day, $j \geq 1$ a mission & \\
$k \in \{1, \ldots, T\}$ & day; index $k = 0$ holds the initial state & \\
$H = [H_1, H_2]$, $T = H_1 + H_2$ & horizon input; a scalar $H$ means $H_1 = H_2 = H$ &
  \ref{sec:input_scalars} \\
days $1, \ldots, H_1$ & transitory phase & \\
days $H_1 + 1, \ldots, T$ & operating phase, repeated without end in the endless
  schedule & \ref{sec:loop} \\
$p \in \{1, \ldots, P\}$ & segment of the \texttt{pwl} form, $P = \texttt{pwl\_points}$ &
  \ref{sec:linear_forms} \\
\midrule
\multicolumn{3}{@{}l}{\emph{Parameters, increments and states}} \\
$\tau_{i\ell}$, $\varepsilon_{i\ell}$, $\rho_{i\ell}$ & failure threshold,
  reliability level, repair efficiency of a cell; a model section may drop the
  subscript $i\ell$ after saying so & \ref{sec:inputs} \\
$C_M$; $C_{R,i\ell}$, $C_{\mathrm{rep},i\ell}$; $C_D$ & costs: depot day; repair,
  replacement; damage penalty & \ref{sec:objective} \\
$\mu^{\inc}_{i\ell jk}$, $v^{\inc}_{i\ell jk}$, $b^{\inc}_{i\ell jk}$,
  $K^{\inc}_{i\ell jk}$, $\beta_{i\ell jk}$ & mean, variance, support width and
  cumulant generating function (CGF) of the increment of cell $(i,\ell)$ when train $i$
  flies mission $j$ on day $k$, after the phase lookup; Gamma rate &
  \ref{sec:input_increments} \\
$\delta\mu_{i\ell k} = \sum_{j=1}^{M} x_{ijk}\, \mu^{\inc}_{i\ell jk}$ & day increment
  of the mean; $\delta v$, $\delta R$ (from $(b^{\inc})^2$) and $\delta K$ likewise; $0$
  on depot and idle days & \\
$\mu_{i\ell k}$, $v_{i\ell k}$, $R_{i\ell k}$, $K_{i\ell k}$ & tracked mean,
  variance, Hoeffding budget, Chernoff quantity & \ref{sec:variables} \\
$g^{\mu}_{i\ell k}$, $g^{v}_{i\ell k}$, $g^{R}_{i\ell k}$, $g^{K}_{i\ell k}$ &
  anchors (law $\ARDone$) & \ref{sec:repair_laws} \\
$\mu_{0,i\ell}$, $v_{0,i\ell}$, $R_{0,i\ell}$, $K_{0,i\ell}$ & initial values &
  \ref{sec:inputs} \\
$\mu^{\new}_{i\ell}$, $v^{\new}_{i\ell}$, $R^{\new}_{i\ell}$, $K^{\new}_{i\ell}$ &
  replacement values & \ref{sec:replacement} \\
$b^{0}_{i\ell}$, $b^{\new}_{i\ell}$ & support widths of the initial and replacement
  damage & \ref{sec:dyn_rainflow} \\
$\ARDinf$, $\ARDone$; $a$; $\psi$ & repair laws; repair factor; repair map &
  \ref{sec:repair_laws} \\
$\xi_k$, $\zeta_k$, $\varphi_k$, $\Phi$, $\mathcal{C}_{i\ell}$ & full tracked state,
  reference trajectory, day map, map of the operating phase, certificate set &
  \ref{sec:encoding}, \ref{sec:loop} \\
\midrule
\multicolumn{3}{@{}l}{\emph{Derived constants}} \\
$M_s$ & natural range of a tracked quantity $s$ (big-M constant) &
  \ref{sec:encoding} \\
$\kappa_G = \Phi^{-1}(1 - \varepsilon)$ & standard normal quantile & \ref{sec:dyn_gaussian} \\
$V_{\max} = \tau^2 / \kappa_G^2$ & variance range, Gaussian and
  Wiener & \ref{sec:dyn_gaussian} \\
$\eta_{i\ell}$; $Z \sim \mathrm{IG}(1, 1/\eta)$; $q$; $\bar\mu = \tau /
  \mathbb{E}[Z \mid Z > q]$ & inverse Gaussian dispersion; unit law; its
  $(1-\varepsilon)$-quantile; mean cap & \ref{sec:dyn_ig} \\
$\beta^*_{i\ell}$; $t = \beta^*\tau$; $A_1(t, \varepsilon)$; $Q(a, x)$ & smallest
  Gamma rate among the terms of the cell with a positive mean; scaled threshold;
  shape cap; regularised upper incomplete gamma function & \ref{sec:dyn_gamma} \\
$\kappa_C = \sqrt{(1-\varepsilon)/\varepsilon}$ & Cantelli coefficient &
  \ref{sec:dyn_rainflow} \\
$\kappa = \ln(1/\varepsilon)$ & exponent constant (Hoeffding, Bernstein) &
  \ref{sec:dyn_rainflow} \\
$b_{i\ell}$, $c_{i\ell} = \kappa\, b_{i\ell}/3$ & Bernstein support constant and margin &
  \ref{sec:dyn_rainflow} \\
$s_{i\ell}$ & Chernoff parameter (input \texttt{s\_chernoff}) & \ref{sec:dyn_rainflow} \\
$V^{\mathrm{C}}$, $V^{\mathrm{B}}$ & variance ranges, rainflow Cantelli and Bernstein &
  \ref{sec:dyn_rainflow} \\
$\mu_r = \texttt{tangent\_ref}\cdot\tau$; $c_2$, $c_1$ & tangent point; constants of a
  quadratic row & \ref{sec:linear_forms} \\
$E$ & evaluation horizon (input \texttt{evaluation\_horizon}) & \ref{sec:objective} \\
\bottomrule
\end{longtable}
\endgroup

% ==== BLOCK 1B ====

% ================================================================
\section{Decision Variables}
\label{sec:variables}
% ================================================================

A variable, and the rows that define it, are built only where a row or an objective
term with a non-zero coefficient uses them. Example: $u$ and its rows are not built
when $C_D = 0$; $z_{i\ell k}$ and its rows are not built at a cell with
$C_{R,i\ell} = 0$. The day index runs over $k = 1, \ldots, T$; index $k = 0$ holds
constants (initial values and starting anchors), not variables.

\begingroup\small
\begin{longtable}{@{}L{2.3cm}L{2.6cm}L{4.2cm}L{5.4cm}@{}}
\toprule
Variable & Shape, type & Built when & Meaning \\
\midrule
\endhead
$x_{ijk}$ & $F \times (M{+}1) \times T$, binary & always &
  1 if train $i$ flies mission $j$ on day $k$; $j = 0$: depot day \\
$x^m_{i\ell k}$ & $F \times L \times T$, binary & always &
  1 if component $\ell$ of train $i$ is repaired on day $k$ \\
$x^r_{i\ell k}$ & $F \times L \times T$, binary & \texttt{allow\_replacement = true} &
  1 if component $\ell$ of train $i$ is replaced on day $k$ \\
$n_{i\ell k}$ & binary & cells with \texttt{encoding = "indicator"} &
  1 on a day without intervention, Eq.~\eqref{eq:a_ind_sum} \\
$\mu_{i\ell k}$ & continuous, $[0, \tau_{i\ell}]$ & every cell &
  tracked mean degradation \\
$v_{i\ell k}$ & continuous, $[0, M_v]$ & cells whose model or bound tracks a variance &
  tracked variance \\
$R_{i\ell k}$ & continuous, $[0, M_R]$ & rainflow cells with the Hoeffding bound &
  tracked Hoeffding budget \\
$K_{i\ell k}$ & continuous, $[0, M_K]$ & rainflow cells with the Chernoff bound &
  tracked Chernoff quantity \\
$g^{\mu}_{i\ell k}$, $g^{v}_{i\ell k}$, $g^{R}_{i\ell k}$, $g^{K}_{i\ell k}$ &
  continuous, range of the quantity & $\ARDone$ cells, one per tracked quantity
  except the Wiener variance & anchors, Eq.~\eqref{eq:a_anchor} \\
$w_{i\ell kp}$ & binary, $p = 1, \ldots, P$ & cells with a quadratic row in the
  \texttt{pwl} form, $P \geq 2$ & segment of the mean, Eq.~\eqref{eq:a_pwl_sum} \\
$z_{i\ell k}$ & continuous, $\geq 0$ & cells with $C_{R,i\ell} > 0$, on days whose cost
  enters the objective & degradation removed by a repair,
  Section~\ref{sec:repair_cost} \\
$u$ & scalar, $\geq 0$ & $C_D > 0$ & damage regularisation,
  Section~\ref{sec:damage_reg} \\
\bottomrule
\end{longtable}
\endgroup

The ranges $M_v$, $M_R$, $M_K$ are the natural ranges of Section~\ref{sec:encoding};
each model section states which quantities it tracks. The output reports the values
of the tracked quantities; they are upper bounds on the reference values
(Section~\ref{sec:encoding}).

% ================================================================
\section{Objective Function}
\label{sec:objective}
% ================================================================

The cost of day $k$ is
\begin{equation}
  c_k = C_M \sum_{i=1}^{F} x_{i0k}
      + \sum_{i=1}^{F}\sum_{\ell=1}^{L} C_{R,i\ell}\, z_{i\ell k}
      + \sum_{i=1}^{F}\sum_{\ell=1}^{L} C_{\mathrm{rep},i\ell}\, x^r_{i\ell k},
  \label{eq:a_daycost}
\end{equation}
where $z_{i\ell k}$ reads $0$ where it is not built, and the last sum is absent when
\texttt{allow\_replacement = false}. The costs of the two phases are
\begin{equation}
  J_{\mathrm{init}} = \sum_{k=1}^{H_1} c_k, \qquad
  J_{\mathrm{op}} = \sum_{k=H_1+1}^{T} c_k .
  \label{eq:a_phases}
\end{equation}
The input \texttt{objective\_mode} (default \texttt{total}) selects the objective:
\begin{equation}
  \min\;
  \begin{cases}
    J_{\mathrm{init}} + J_{\mathrm{op}} + C_D\, u & \texttt{total}, \\[2pt]
    J_{\mathrm{op}} / H_2 + C_D\, u & \texttt{operating\_average}, \\[2pt]
    J_{\mathrm{init}} + \dfrac{E - H_1}{H_2}\, J_{\mathrm{op}} + C_D\, u
      & \texttt{evaluation\_total},
  \end{cases}
  \label{eq:a_objective}
\end{equation}
where $E$ is the input \texttt{evaluation\_horizon}, required in the mode
\texttt{evaluation\_total}, with $E > H_1$. The term $C_D\,u$ is present only when
$C_D > 0$ (Section~\ref{sec:damage_reg}).

\textbf{Cost terms.}
\begin{itemize}[nosep,leftmargin=1.5em]
  \item $C_M \geq 0$, a scalar, is paid once per train per depot day ($x_{i0k} = 1$),
        whatever the number of components repaired or replaced on that day. It
        prices the visit to the depot, for example the train out of service and the
        set-up. An idle day (Section~\ref{sec:assignment}) costs nothing, so the
        number of depot days is a costed decision.
  \item $C_{R,i\ell} \geq 0$ is paid per unit of mean degradation removed by a
        repair (Section~\ref{sec:repair_cost}).
  \item $C_{\mathrm{rep},i\ell} \geq 0$ is paid per replacement. The input is
        required; it has no default.
  \item $C_D \geq 0$ weights the damage penalty $u$
        (Section~\ref{sec:damage_reg}). The term is added once in every mode,
        because $u$ is the maximum over the whole horizon.
\end{itemize}
The shapes of the cost inputs are given in Section~\ref{sec:inputs}.

\textbf{Meaning of the modes.}
\begin{itemize}[nosep,leftmargin=1.5em]
  \item \texttt{total}: the cost of the $T$ modelled days.
  \item \texttt{operating\_average}: the cost per day of the operating phase. It is
        the long-run cost per day of the endless schedule
        (Section~\ref{sec:loop}), in which the cost of the transitory phase spreads
        over infinitely many days. It compares solves with different $H_2$
        (Section~\ref{sec:sweeps}). The transitory costs do not enter the
        objective in this mode: the solver chooses the transitory schedule for
        feasibility only.
  \item \texttt{evaluation\_total}: the projected cost over $E$ days: the transitory
        phase once, then $(E - H_1)/H_2$ repetitions of the operating phase, a
        fraction of a repetition counted pro rata. It also prices the transitory
        phase, so it compares solves with different $H_1$.
\end{itemize}
Example: $H_1 = 5$, $H_2 = 10$, $E = 365$. Then \texttt{evaluation\_total} weights
the operating cost by $36$, and \texttt{operating\_average} by $1/10$.

The objective value is computed from the tracked values. It is therefore an upper
bound on the cost of the schedule computed from the reference values
(Section~\ref{sec:encoding}).

% ================================================================
\section{Constraints}
\label{sec:constraints}
% ================================================================

The constraints have three groups. The fleet rows (assignment and depot days) do not
depend on the degradation models. The common rules of the dynamics (repair laws,
encoding, replacement, linear forms of the reliability rows) hold for every model;
the five model sections apply them and add each model's tracked quantities, day
maps and reliability row. The repair cost, the damage regularisation and the loop
constraint close the problem. All rows hold for every day $k = 1, \ldots, T$ unless a
section says otherwise.

\subsection{Assignment Constraints}
\label{sec:assignment}

Each train performs at most one activity per day:
\begin{equation}
  \sum_{j=0}^{M} x_{ijk} \leq 1 \qquad \forall\, i, k .
  \label{eq:a_one_activity}
\end{equation}
Each mission is flown by exactly one train every day:
\begin{equation}
  \sum_{i=1}^{F} x_{ijk} = 1 \qquad \forall\, j \in \{1, \ldots, M\},\; k .
  \label{eq:a_coverage}
\end{equation}
Together these rows put exactly $M$ trains on missions every day. The other $F - M$
trains are in the depot or idle. No other fleet-availability row is needed.

\emph{Idle days.} A train is idle on day $k$ when its assignment row is all zeros:
$x_{ijk} = 0$ for every $j \in \{0, 1, \ldots, M\}$. An idle train flies no mission
and is not in the depot. Idle days and depot days add no increment, because every
day aggregate $\delta\mu_{i\ell k}$, $\delta v_{i\ell k}$, $\delta R_{i\ell k}$,
$\delta K_{i\ell k}$ sums over missions $j \geq 1$ only. A depot day must be chosen
explicitly ($x_{i0k} = 1$), and repairs and replacements need one
(Section~\ref{sec:depot}). Because Eq.~\eqref{eq:a_one_activity} is an inequality,
the choice between a depot day and an idle day is free, and the term $C_M \sum
x_{i0k}$ of the objective prices it.

\subsection{Depot-Day Compatibility}
\label{sec:depot}

A component may be repaired or replaced only on a depot day of its train, and at most
one of the two on the same day:
\begin{equation}
  x^m_{i\ell k} + x^r_{i\ell k} \leq x_{i0k} \qquad \forall\, i, \ell, k .
  \label{eq:a_depot}
\end{equation}
With \texttt{allow\_replacement = false} the row reads $x^m_{i\ell k} \leq x_{i0k}$.
Several components of one train may receive interventions on the same depot day. A
depot day without an intervention on a component leaves that component unchanged.

\emph{Remark (one row instead of three).} Eq.~\eqref{eq:a_depot} has the same binary
solutions as the three rows $x^m \leq x_{i0k}$, $x^r \leq x_{i0k}$, $x^m + x^r \leq 1$.
Its linear relaxation is tighter: the point $x_{i0k} = x^m = x^r = 0.5$ satisfies the
three rows but not Eq.~\eqref{eq:a_depot}. With $0 \leq x^m, x^r$ and $x_{i0k} \leq 1$,
Eq.~\eqref{eq:a_depot} describes the convex hull of the four binary solutions
$(x_{i0k}, x^m, x^r) \in \{(0,0,0), (1,0,0), (1,1,0), (1,0,1)\}$.

\subsection{Repair Laws}
\label{sec:repair_laws}

An imperfect repair (in short, a repair) of a component on a depot day removes part
of its degradation. The repair efficiency $\rho_{i\ell} \in (0, 1]$ is the fraction
removed. Let $D^-$ be the degradation level before the repair (at the end of day
$k-1$; a depot day adds no increment) and $D^+$ the level after it (day $k$). The
input \texttt{repair\_model} selects one of two laws per cell; the default is
\texttt{ardinf}.

\begin{definition}[Repair laws]
\label{def:a_laws}
Let $\rho = \rho_{i\ell}$. The two repair laws are:
\begin{itemize}[nosep,leftmargin=1.5em]
  \item $\ARDinf$ (input value \texttt{ardinf}) removes the fraction $\rho$ of the
        whole level:
        \begin{equation}
          D^+ = (1-\rho)\, D^- .
          \label{eq:a_ardinf}
        \end{equation}
  \item $\ARDone$ (input value \texttt{ard1}) removes the fraction $\rho$ of the
        degradation accumulated since the last intervention (repair or
        replacement):
        \begin{equation}
          D^+ = D^- - \rho\,\bigl(D^- - D^{\mathrm{last}}\bigr),
          \label{eq:a_ard1}
        \end{equation}
        where $D^{\mathrm{last}}$ is the level just after the last intervention.
        Before the first intervention of the horizon, $D^{\mathrm{last}}$ is the
        level just after the last intervention before the horizon.
\end{itemize}
\end{definition}

\emph{Names and sources.} $\ARDone$, ``arithmetic reduction of deterioration of
order 1'', is the law of \cite[$\S$3.1, eq.~(4)]{mercier2019}. That source measures the
removed degradation from the level just after the last maintenance, as
Eq.~\eqref{eq:a_ard1} does. The name $\ARDinf$ is an analogy, not a citation.
Doyen and Gaudoin define the classes ARA$_\infty$ \cite[eq.~(20), p.~11]{doyen2004}
and ARA$_1$ \cite[eq.~(22), p.~12]{doyen2004}, which reduce a virtual age, not a
degradation level. ARA$_\infty$ reduces the whole virtual age; ARA$_1$ reduces only
the age accumulated since the previous repair. The subscript is the memory of the
class: the largest number of previous failure times that enter the failure
intensity \cite[p.~8]{doyen2004}. $\ARDinf$ applies the rule of ARA$_\infty$ to the
whole degradation level. No law of this specification acts on a virtual age.

\textbf{Repair maps of the tracked quantities.} A model tracks moments or other
summaries of $D$, not $D$ itself. For a tracked quantity $s$ with anchor $g$, a
repair on day $k$ applies the repair map
\begin{equation}
  \psi(s_{k-1}, g_{k-1}) =
  \begin{cases}
    a\, s_{k-1} & \ARDinf, \\
    a\, s_{k-1} + (1-a)\, g_{k-1} & \ARDone,
  \end{cases}
  \label{eq:a_psi}
\end{equation}
with a repair factor $a \in [0, 1]$. For the mean, $a = 1 - \rho_{i\ell}$, because
expectation is linear: Eq.~\eqref{eq:a_ard1} gives $\mathbb{E}[D^+] =
(1-\rho)\,\mathbb{E}[D^-] + \rho\,\mathbb{E}[D^{\mathrm{last}}]$. For a variance and
for $R$ the factor is $(1-\rho_{i\ell})^2$, and for $K$ it is $1-\rho_{i\ell}$; the
model sections derive each factor. The Wiener variance has the factor $a = 1$
under both laws: a repair acts on the Wiener mean only and leaves the variance
unchanged.

\textbf{Anchors.} Under $\ARDone$ each tracked quantity has an anchor $g$: the value
of the quantity just after the last intervention. The Wiener variance is the
exception: with $a = 1$ its map does not use an anchor, and none is built. The
anchor recursion is
\begin{equation}
  g_0 = \text{starting anchor}, \qquad
  g_k = \begin{cases}
    s_k & \text{on a repair or replacement day}, \\
    g_{k-1} & \text{otherwise},
  \end{cases}
  \label{eq:a_anchor}
\end{equation}
built by the rows Eq.~\eqref{eq:a_bigm_anchor} or Eq.~\eqref{eq:a_ind_anchor}. The
starting anchors $g^{\mu}_{i\ell 0}$, $g^{v}_{i\ell 0}$, $g^{R}_{i\ell 0}$,
$g^{K}_{i\ell 0}$ come from the optional inputs \texttt{anchor\_mu\_0},
\texttt{anchor\_v\_0}, \texttt{anchor\_R\_0}, \texttt{anchor\_K\_0}. The default of
each is the initial value of its quantity. Each must satisfy $0 \leq g_0 \leq s_0$
(Section~\ref{sec:consistency_checks}). Example: $\mu_{k-1} = 0.6$, $g_{k-1} = 0.2$,
$\rho = 0.5$. A repair gives $\mu_k = 0.6 - 0.5\,(0.6 - 0.2) = 0.4$ and $g_k = 0.4$.
On the next day the gap $\mu - g$ is $0$, and it grows again with each mission.

\emph{Remark (default of the starting anchor).} Damage only grows between
interventions. So the true level after the last intervention before the horizon lies
between the level of a new part and the initial level. A larger anchor gives a
weaker repair and a higher state, because $\psi$ is non-decreasing in $g$. The
default $g_0 = s_0$ therefore never overstates what a repair achieves; the input
gives the exact value when the maintenance history is known. Example: $\mu_0 =
0.5$, $\rho = 0.5$, and a repair on day~1 before any mission. With the starting
anchor $0$ the mean drops to $0.25$; with the default $\mu_0$ it stays at $0.5$, and
the repair removes only damage accumulated within the horizon.

\textbf{Laws per model.} Every model accepts both laws. The table below gives the
reason for each model.

\begingroup\small
\begin{longtable}{@{}lllp{7.4cm}@{}}
\toprule
Model & $\ARDinf$ & $\ARDone$ & Reason or note \\
\midrule
Gaussian & yes & yes & both laws keep the state Gaussian, with exact moments \\
Inverse Gaussian & yes & yes & the cap holds for any repair weights and any
  dependence between the terms (Section~\ref{sec:dyn_ig}) \\
Wiener & yes & yes & on the mean only; the variance changes only by replacement \\
Gamma & yes & yes & Section~\ref{sec:dyn_gamma} \\
Rainflow, every bound & yes & yes & Chernoff under $\ARDone$ uses a conservative
  recursion for $K$ (Section~\ref{sec:dyn_rainflow}) \\
\bottomrule
\end{longtable}
\endgroup

\subsection{Encoding of the Dynamics}
\label{sec:encoding}

Each day, every tracked quantity of a cell follows one rule: a mission adds an
increment, a repair applies the repair map of Section~\ref{sec:repair_laws}, and a
replacement resets the quantity (Section~\ref{sec:replacement}). This subsection
states these rules once, for every model and both repair laws. The model sections
give, for each tracked quantity, its increment, its repair factor and its natural
range.

\textbf{One-sided rows.} Every dynamics row is an inequality ($\geq$), not an
equality. The solver may choose a value above the result of the rule, never a
value below it. So a tracked quantity is an upper bound on the value that the rule
gives. Example: with $\mu_{i\ell,k-1} = 0.3$ and a mission increment of $0.2$, the
rule gives $0.5$; the rows accept $\mu_{i\ell k} = 0.5$ and also $\mu_{i\ell k} =
0.6$. Replacement rows and anchor updates are equalities.

\textbf{Generic rows.} Let $s$ be one tracked quantity of a cell (a mean, a
variance, $R$ or $K$), with day increment $\delta s_k$, natural range $M_s$,
replacement value $s^{\new}$, anchor $g$ (law $\ARDone$ only) and repair map
$\psi$ (Section~\ref{sec:repair_laws}). Write $x^m_k$, $x^r_k$ for the binaries of
the cell and $e_k = x^m_k + x^r_k$. The input \texttt{encoding} selects one of two
forms per cell. Both forms are built for $k = 1, \ldots, T$; index $k = 0$ holds the
initial value and the starting anchor.

\emph{Big-M form} (\texttt{encoding = "bigm"}, the default):
\begin{align}
  &\text{no intervention:} & s_k &\geq s_{k-1} + \delta s_k - M_s\,(x^m_k + x^r_k),
    \label{eq:a_bigm_carry} \\
  &\text{repair:} & s_k &\geq \psi(s_{k-1}, g_{k-1}) - M_s\,(1 - x^m_k),
    \label{eq:a_bigm_repair} \\
  &\text{replacement:} & |s_k - s^{\new}| &\leq M_s\,(1 - x^r_k),
    \label{eq:a_bigm_repl} \\
  &\text{anchor } (\ARDone)\text{:} & |g_k - s_k| &\leq M_s\,(1 - e_k), \qquad
    |g_k - g_{k-1}| \leq M_s\, e_k .
    \label{eq:a_bigm_anchor}
\end{align}
Each row with an absolute value is built as two linear rows.

\emph{Indicator form} (\texttt{encoding = "indicator"}). One extra binary
$n_{i\ell k}$ per cell and day marks a day without intervention:
\begin{equation}
  n_{i\ell k} + x^m_{i\ell k} + x^r_{i\ell k} = 1 .
  \label{eq:a_ind_sum}
\end{equation}
The rows are Gurobi indicator constraints:
\begin{align}
  n_k = 1 &\;\Rightarrow\; s_k \geq s_{k-1} + \delta s_k, &
  x^m_k = 1 &\;\Rightarrow\; s_k \geq \psi(s_{k-1}, g_{k-1}), &
  x^r_k = 1 &\;\Rightarrow\; s_k = s^{\new}, \label{eq:a_ind_state} \\
  n_k = 1 &\;\Rightarrow\; g_k = g_{k-1}, &
  x^m_k = 1 &\;\Rightarrow\; g_k = s_k, &
  x^r_k = 1 &\;\Rightarrow\; g_k = s_k. \label{eq:a_ind_anchor}
\end{align}
The binary $n_{i\ell k}$ is built only at cells in this form.

\emph{Common rules for both forms.}
\begin{itemize}[nosep,leftmargin=1.5em]
  \item The anchor rows are built only under $\ARDone$, one anchor per tracked
        quantity, except the Wiener variance.
  \item The Wiener variance takes the repair map $\psi = s_{k-1}$ (factor $a = 1$)
        under both laws, and has no anchor. A repair day adds no increment, so in the
        big-M form its no-intervention and repair rows together equal the single row
        $s_k \geq s_{k-1} + \delta s_k - M_s\, x^r_k$ (Eq.~\eqref{eq:w_vrow}), which
        may replace them. The variance is never left free below on a repair day.
  \item With \texttt{allow\_replacement = false} the binaries $x^r$ are not built:
        read $x^r_k = 0$ in every row, and drop the replacement rows.
  \item At binary points the two forms have the same feasible set. Every statement
        below holds for both.
\end{itemize}

\textbf{Natural ranges.} The big-M constant $M_s$ is the natural range of $s$: an
upper bound of $s$ at every point that satisfies the reliability row of the cell.
It is $\tau_{i\ell}$ for a mean and for its anchor, $V_{\max}$ for a Gaussian or
Wiener variance, $V^{\mathrm{C}}$ or $V^{\mathrm{B}}$ for a rainflow variance
(Cantelli or Bernstein), $2\tau^2/\kappa$ for $R$ and $s_{i\ell}\tau - \kappa$ for
$K$ (Section~\ref{sec:dyn_rainflow}). An anchor has the range of its quantity.

\emph{Remark (validity of the big-M constants).} Every tracked quantity and every
anchor lies in $[0, M_s]$. The initial values, the starting anchors and the
replacement values lie in $[0, M_s]$ too, by the input checks of
Section~\ref{sec:consistency_checks}. An intervention day is a depot day
(Section~\ref{sec:depot}), so $\delta s_k = 0$ on it. A relaxed row then cuts no
point: Eq.~\eqref{eq:a_bigm_carry} with $e_k = 1$ reads
$s_k \geq s_{k-1} - M_s$, whose right-hand side is at most $0$;
Eq.~\eqref{eq:a_bigm_repair} with $x^m_k = 0$ has a right-hand side
$\psi - M_s \leq 0$, because $\psi$ is a convex combination of values in
$[0, M_s]$; Eq.~\eqref{eq:a_bigm_repl} with $x^r_k = 0$ allows
$s_k \in [s^{\new} - M_s,\, s^{\new} + M_s] \supseteq [0, M_s]$; the relaxed anchor
rows allow every difference in $[-M_s, M_s]$.

\emph{Remark (choice of the default form).} Gurobi accepts a binary as integral
within its integrality tolerance \texttt{IntFeasTol} (default $10^{-5}$). A big-M
row can therefore leak by at most $10^{-5} M_s$. Because $M_s$ is the quantity's own
range, this is a relative error of $10^{-5}$. Indicator constraints have no such
leak, and Gurobi usually needs more time to solve them. Which form solves this
model faster has not been measured.

\subsubsection{Safety of the one-sided rows}

Fix a schedule, that is, values of every binary, and one cell. The following
objects are used here and in Section~\ref{sec:loop}.
\begin{itemize}[nosep,leftmargin=1.5em]
  \item The \emph{full tracked state} $\xi_k$ collects the values on day $k$ of every
        tracked quantity of the cell and of every anchor. $\xi_0$ holds the initial
        values and the starting anchors.
  \item The \emph{day map} $\varphi_k$ gives, for each coordinate, the value that the
        rule of day $k$ assigns: $s + \delta s_k$ and $g$ on a day without
        intervention; $\psi(s, g)$ for both $s$ and $g$ on a repair day; $s^{\new}$
        for both on a replacement day.
  \item The \emph{reference trajectory} is $\zeta_0 = \xi_0$,
        $\zeta_k = \varphi_k(\zeta_{k-1})$: the rows with every inequality replaced
        by an equality.
  \item The \emph{true statistics} $\theta_k$ are the true values on day $k$ of the
        tracked quantities (the mean, the variance, $R$ or the CGF of the damage
        $D_{i\ell k}$) and of their anchors (their values just after the last
        intervention).
  \item The \emph{certificate set} $\mathcal{C}_{i\ell}$ is the set of values of the
        tracked quantities at which the reliability row of the cell holds in its
        exact form, side rows included. Anchors do not enter it.
\end{itemize}

Every model section states three facts, labelled (M1)--(M3) there:
\begin{enumerate}[nosep,leftmargin=2em,label=(M\arabic*)]
  \item Each repair factor $a$ lies in $[0,1]$ and each increment
        $\delta s_k$ is non-negative.
  \item On day $0$ the true statistics are at most the initial values and starting
        anchors, $\theta_0 \leq \xi_0$. On every later day the true statistic is at
        most the day map applied to the true statistics of the day before:
        $\theta_k \leq \varphi_k(\theta_{k-1})$. Every map holds with equality
        except the map of the Chernoff quantity $K$, which is an upper bound
        (Section~\ref{sec:dyn_rainflow}).
  \item If values $s \in \mathcal{C}_{i\ell}$ are at least the true statistics of
        day $k$, then $\Pr(D_{i\ell k} > \tau_{i\ell}) \leq \varepsilon_{i\ell}$.
\end{enumerate}
The set $\mathcal{C}_{i\ell}$ is closed, and it is a down-set: if $s \in
\mathcal{C}_{i\ell}$ and $0 \leq s' \leq s$ componentwise, then $s' \in
\mathcal{C}_{i\ell}$. The reason is that every reliability row, written in
chance-bound form, reads ``a continuous function of the tracked quantities,
non-decreasing in each of them, is at most a constant'':

\begingroup\small
\begin{longtable}{@{}lll@{}}
\toprule
Model or bound & Row in chance-bound form & Section \\
\midrule
Gaussian, Wiener & $\mu + \Phi^{-1}(1-\varepsilon)\sqrt{v} \leq \tau$ &
  \ref{sec:dyn_gaussian}, \ref{sec:dyn_wiener} \\
Inverse Gaussian; Gamma & $\mu \leq \bar\mu$; $\mu \leq A_1(\beta^*\tau,\varepsilon)/\beta^*$ &
  \ref{sec:dyn_ig}, \ref{sec:dyn_gamma} \\
Rainflow, Markov & $\mu \leq \varepsilon\,\tau$ & \ref{sec:dyn_rainflow} \\
Rainflow, Cantelli & $\mu + \kappa_C\sqrt{v} \leq \tau$ & \ref{sec:dyn_rainflow} \\
Rainflow, Hoeffding & $\mu + \sqrt{\kappa R/2} \leq \tau$ & \ref{sec:dyn_rainflow} \\
Rainflow, Bernstein & $\mu + c + \sqrt{c^2 + 2\kappa v} \leq \tau$ & \ref{sec:dyn_rainflow} \\
Rainflow, Chernoff & $K \leq s_{i\ell}\,\tau + \ln\varepsilon$ & \ref{sec:dyn_rainflow} \\
\bottomrule
\end{longtable}
\endgroup

\begin{lemma}[Monotone day maps]
\label{lem:a_monotone}
Each day map $\varphi_k$ is non-decreasing in each coordinate of the full tracked
state, anchors included.
\end{lemma}
\begin{proof}
Let $\xi \leq \xi'$ componentwise. Each coordinate of $\varphi_k$ has one of the
forms $s + \delta s_k$, $g$, $s$, $a\,s$, $a\,s + (1-a)\,g$ or $s^{\new}$. By (M1),
$a \geq 0$ and $1 - a \geq 0$. Therefore each coordinate of $\varphi_k(\xi)$ is at
most the same coordinate of $\varphi_k(\xi')$.
\end{proof}

\begin{proposition}[Domination]
\label{prop:a_domination}
Every feasible solution satisfies $\xi_k \geq \varphi_k(\xi_{k-1})$ and
$\theta_k \leq \zeta_k \leq \xi_k$ for $k = 1, \ldots, T$.
\end{proposition}
\begin{proof}
At binary points the rows of day $k$ read: $s_k \geq \varphi_k(\xi_{k-1})_s$ for
every tracked quantity $s$, with equality on a replacement day; $g_k = s_k \geq
\varphi_k(\xi_{k-1})_g$ on an intervention day; $g_k = g_{k-1} =
\varphi_k(\xi_{k-1})_g$ on any other day. So $\xi_k \geq \varphi_k(\xi_{k-1})$.
Let $\theta_{k-1} \leq \zeta_{k-1} \leq \xi_{k-1}$ (true for $k = 1$ by (M2), since
$\zeta_0 = \xi_0$). Then, by (M2) and Lemma~\ref{lem:a_monotone},
\[
  \theta_k \leq \varphi_k(\theta_{k-1}) \leq \varphi_k(\zeta_{k-1}) = \zeta_k
  \leq \varphi_k(\xi_{k-1}) \leq \xi_k .
\]
Therefore $\theta_k \leq \zeta_k \leq \xi_k$ for every $k$ by induction.
\end{proof}

\emph{Consequence on the horizon.} On every day $k = 1, \ldots, T$ the tracked
values satisfy the reliability row in the form chosen for the cell. Every form lies
inside $\mathcal{C}_{i\ell}$ (Section~\ref{sec:linear_forms}), so $\xi_k \in
\mathcal{C}_{i\ell}$. By Proposition~\ref{prop:a_domination}, $\theta_k \leq \xi_k$,
so (M3) gives $\Pr(D_{i\ell k} > \tau_{i\ell}) \leq \varepsilon_{i\ell}$.
Theorem~\ref{thm:a_loop} (Section~\ref{sec:loop}) extends this to every day of
every later repetition of the operating phase.

\textbf{What the one-sided rows accept.} Call a schedule \emph{safe} for a cell when
the reference trajectory of the endless schedule (the transitory phase once, then
the operating phase repeated without end, Section~\ref{sec:loop}) lies in
$\mathcal{C}_{i\ell}$ and in the natural ranges on every day. The one-sided rows
relate to this property as follows.
\begin{itemize}[nosep,leftmargin=1.5em]
  \item \emph{They never accept an unsafe schedule} (Theorem~\ref{thm:a_loop}).
  \item \emph{They accept every schedule that equality rows accept}, and more.
        Equality rows force $\xi_k = \zeta_k$; Proposition~\ref{prop:a_accept}~(a)
        shows that such a schedule satisfies the one-sided rows too.
        The first example below shows a safe schedule that only the one-sided
        rows accept.
  \item \emph{They may reject a safe schedule}, as equality rows may, but only when
        tracked quantities move in opposite directions over the first repetition
        (Proposition~\ref{prop:a_accept}~(b) and the second example below).
\end{itemize}

\begin{proposition}[Acceptance]
\label{prop:a_accept}
Fix a schedule and a cell, use the exact reliability form, and let the schedule be
safe for the cell. Let $\bar\zeta = \zeta_{H_1}$ and let $\Phi$ be the composition
of the day maps of the operating phase (Section~\ref{sec:loop}).
\begin{enumerate}[nosep,label=(\alph*)]
  \item If $\Phi(\bar\zeta) \leq \bar\zeta$, the rows of the cell, the loop row
        included, accept the schedule with $\xi_k = \zeta_k$ for every $k$.
  \item If $\Phi(\bar\zeta) \geq \bar\zeta$, the cell has no anchor (law
        $\ARDinf$), and the cell is not replaced on day $H_1$, the rows of the cell
        accept the schedule.
\end{enumerate}
\end{proposition}
\begin{proof}
(a) The values $\xi_k = \zeta_k$ satisfy every row of day $k$ with equality at
binary points. They lie in $\mathcal{C}_{i\ell}$ because the schedule is safe. The
loop row reads $\zeta_T = \Phi(\bar\zeta) \leq \bar\zeta = \zeta_{H_1}$.

(b) Let $y_n = \Phi^n(\bar\zeta)$, the reference state at the end of the $n$-th
repetition. By Lemma~\ref{lem:a_monotone} and $\Phi(\bar\zeta) \geq \bar\zeta$,
$y_{n+1} = \Phi(y_n) \geq \Phi(y_{n-1}) = y_n$. The sequence is non-decreasing and
lies in the natural ranges, so it converges to some $y^*$ with $y^* \geq
\bar\zeta$. Each day map is affine or constant, so $\Phi$ and its prefixes
$\Phi_h$ are continuous: $\Phi(y^*) = \lim_n y_{n+1} = y^*$ and
$\Phi_h(y^*) = \lim_n \Phi_h(y_n)$. Each $\Phi_h(y_n)$ is a reference value of a
later day, so it lies in the closed set $\mathcal{C}_{i\ell}$; hence
$\Phi_h(y^*) \in \mathcal{C}_{i\ell}$ for $h = 1, \ldots, H_2$, and
$y^* = \Phi_{H_2}(y^*) \in \mathcal{C}_{i\ell}$. Take $\xi_k =
\zeta_k$ for $k < H_1$, $\xi_{H_1} = y^*$ and $\xi_{H_1+h} = \Phi_h(y^*)$. Day $H_1$
is not a replacement day and the cell has no anchor, so $\xi_{H_1} = y^* \geq
\varphi_{H_1}(\xi_{H_1-1})$ satisfies the rows of day $H_1$. The later days satisfy
their rows with equality, and the loop row reads $\xi_T = \Phi(y^*) = y^* =
\xi_{H_1}$.
\end{proof}

A scalar state is always comparable with its image. So a cell with one tracked
quantity under $\ARDinf$, not replaced on day $H_1$, accepts every safe schedule.
Under $\ARDone$ the anchor is a second tracked quantity, and its value on day $H_1$
is fixed by the anchor rows; part (b) does not apply, and acceptance of every safe
schedule is not proven.

\emph{Example (one quantity, accepted only by the one-sided rows).}
One tracked mean, law $\ARDinf$, $\rho = 0.5$. The operating phase is a mission
($\delta\mu = 0.2$) followed by a repair, so $\Phi(y) = 0.5\,(y + 0.2) = 0.5\,y +
0.1$, and $\Phi(y) \leq y$ if and only if $y \geq 0.2$. Let $\bar\zeta = 0.1$. Then
$\Phi(0.1) = 0.15 > 0.1$, so equality rows reject the schedule. The one-sided rows
accept it with $\xi_{H_1} = 0.2$: the tracked cycle is $0.2, 0.4, 0.2$, and the
reliability row is checked at $0.4$. The reference trajectory is $0.1, 0.3, 0.15,
0.35, 0.175, \ldots$; its start-of-cycle values rise to $0.2$ and its peaks rise to
$0.4$. So the schedule is safe whenever $0.4 \in \mathcal{C}_{i\ell}$. The price
is in the reported cost: each repair is charged $\rho \cdot 0.4 = 0.2$, while the
reference removal rises from $0.15$ toward $0.2$.

\emph{Example (two quantities that move in opposite directions, rejected).}
Cantelli bound, $\varepsilon = 0.1$ ($\kappa_C^2 = 9$), $\tau = 1$, $\rho = 0.5$,
law $\ARDinf$. The operating phase is a mission, which adds $0.2$ to the mean and
$0.009$ to the variance, followed by a repair. So $\Phi(\mu, v) = (0.5\,\mu + 0.1,\;
0.25\,v + 0.00225)$, with fixed point $(0.2, 0.003)$. Let $\bar\zeta = (0.5, 0)$.
\begin{itemize}[nosep,leftmargin=1.5em]
  \item Reference: over the repetitions the mean falls to $0.2$ and the variance
        rises to $0.003$. The first peak is $(0.7, 0.009)$, with $9 \cdot 0.009 =
        0.081 \leq 0.3^2 = 0.09$. Every later peak has a mean of at most $0.55$ and a
        variance below $0.012$, with $9 \cdot 0.012 = 0.108 \leq 0.45^2 = 0.2025$.
        The schedule is safe.
  \item One-sided rows: the loop row needs $\xi_{H_1} = y \geq (0.5, 0)$ with
        $\Phi(y) \leq y$, so $y_v \geq 0.003$. The tracked peak is then at least
        $(0.7, 0.012)$, where $9 \cdot 0.012 = 0.108 > 0.09$. Rejected.
  \item Equality rows: $\Phi(\bar\zeta) = (0.35, 0.00225)$ has a larger variance
        than $\bar\zeta$. Rejected as well.
\end{itemize}

\textbf{Meaning of the reported values.} The output reports the tracked values
$\xi_k$, not the reference values. By Proposition~\ref{prop:a_domination} each
reported state and anchor is an upper bound on its reference value, and so on its
true statistic. The costs
computed from them are upper bounds too: the repair cost
(Proposition~\ref{prop:a_repair_cost}) and the damage penalty
(Section~\ref{sec:damage_reg}).

\emph{Remark (slack in the reported values).} The minimisation removes slack only
where slack costs something: a larger mean raises the damage penalty $C_D\,u$ where
$u$ is binding, and raises the repair cost $C_R\,z$ of every later repair. A larger
variance, $R$ or $K$ tightens the reliability row and so restricts the schedule.
Elsewhere the solver may return any value above the reference value. With $C_D = 0$
the damage penalty is absent (Section~\ref{sec:damage_reg}), and the reported states
may carry more slack. Slack never affects safety (Theorem~\ref{thm:a_loop}).

\subsection{Replacement}
\label{sec:replacement}

A replacement installs a new part on a depot day. The input
\texttt{allow\_replacement} (default \texttt{true}) makes replacements available. When
it is \texttt{false}, the binaries $x^r_{i\ell k}$, the replacement rows and the
replacement cost are not built.

\textbf{State after a replacement.} On a replacement day every tracked quantity of
the cell takes its replacement value, and every anchor takes the same value
(Eq.~\eqref{eq:a_anchor}). The values come from the input, for every model:
\begin{itemize}[nosep,leftmargin=1.5em]
  \item mean $\mu^{\new}_{i\ell}$, input \texttt{replacement\_mu}, default $0$;
  \item variance $v^{\new}_{i\ell}$, input \texttt{replacement\_v}, default $0$, where
        the model tracks a variance;
  \item $R^{\new}_{i\ell}$ and $K^{\new}_{i\ell}$ for the rainflow Hoeffding and
        Chernoff bounds, with their defaults and their derivation in
        Section~\ref{sec:dyn_rainflow}.
\end{itemize}
A replacement value describes the damage of the new part: a known damage when its
variance is $0$, a random one otherwise. The rows are the equalities
Eq.~\eqref{eq:a_bigm_repl} or the indicator rows Eq.~\eqref{eq:a_ind_state}. The
input checks of Section~\ref{sec:consistency_checks} require each replacement value
to lie in the natural range of its quantity, and the replacement state to satisfy the
reliability row of the cell; otherwise no replacement of that cell would be
feasible.

\textbf{Shortcut for a new part.} The input \texttt{replacement\_as\_new} (default
\texttt{false}) states that a replaced part is new, with no damage: every
replacement value of the cell is $0$. A non-zero replacement input together with
\texttt{replacement\_as\_new = true} is an input error
(Section~\ref{sec:consistency_checks}).

\textbf{Cost.} A replacement costs $C_{\mathrm{rep},i\ell}$ and nothing else: the
repair cost $z_{i\ell k}$ is $0$ on a replacement day
(Section~\ref{sec:repair_cost}). A cell may be replaced on any depot day of its
train, any number of times. The loop row can require at least one replacement in
the operating phase: for every $\ARDone$ cell with $\rho < 1$ that accumulates damage
(Proposition~\ref{prop:a_ard1_replacement}), and for the Wiener variance
(Section~\ref{sec:dyn_wiener}).

\subsection{Linear Forms of the Reliability Rows}
\label{sec:linear_forms}

Five reliability rows are quadratic: the Gaussian and Wiener rows, and the rainflow
Cantelli, Hoeffding and Bernstein rows. Each has the form
\begin{equation}
  \chi_{i\ell k} \;\leq\; f(\mu_{i\ell k}),
  \qquad f(\mu) = c_2\,(\tau - \mu)^2 + c_1\,(\tau - \mu),
  \qquad \mu_{i\ell k} \leq \bar\mu_{\mathrm{side}},
  \label{eq:a_quad}
\end{equation}
with the constants of the table below; the model sections derive each row. The
third relation is the side row; it
selects the branch of the parabola on which the row is a valid certificate. Since
$c_2 > 0$, $f$ is convex in $\mu$, with $f'(\mu) = -2c_2(\tau - \mu) - c_1$.

\begingroup\small
\begin{longtable}{@{}lllll@{}}
\toprule
Row & $\chi$ & $c_2$ & $c_1$ & $\bar\mu_{\mathrm{side}}$ \\
\midrule
Gaussian, Wiener & $v$ & $1/(\Phi^{-1}(1-\varepsilon))^2$ & $0$ & $\tau$ \\
Rainflow, Cantelli & $v$ & $1/\kappa_C^2 = \varepsilon/(1-\varepsilon)$ & $0$ & $\tau$ \\
Rainflow, Hoeffding & $R$ & $2/\kappa$ & $0$ & $\tau$ \\
Rainflow, Bernstein & $v$ & $1/(2\kappa)$ & $-b_{i\ell}/3$ & $\tau - c_{i\ell}$ \\
\bottomrule
\end{longtable}
\endgroup

The input \texttt{reliability\_impl} selects one of three forms per cell. Every form
keeps the side row.
\begin{itemize}[leftmargin=1.5em]
  \item \texttt{exact}: Eq.~\eqref{eq:a_quad} as written. The feasible set
        $\{(\mu, \chi) : \chi \leq f(\mu)\}$ lies below a convex function, so it is not
        convex: the row is a nonconvex quadratic constraint. Gurobi's parameter
        \texttt{NonConvex} is set to $2$ whenever at least one cell has a quadratic
        row in this form.
  \item \texttt{tangent}: one linear row per day, the tangent of $f$ at the point
        $\mu_r = \texttt{tangent\_ref}_{i\ell}\cdot\tau_{i\ell}$:
        \begin{equation}
          \chi_{i\ell k} \;\leq\; f(\mu_r) + f'(\mu_r)\,(\mu_{i\ell k} - \mu_r).
          \label{eq:a_tangent}
        \end{equation}
        The input \texttt{tangent\_ref} is per cell, in $[0, 1)$, default $0.5$.
        At \texttt{tangent\_ref} $= 0$ the Gaussian row becomes
        $(\Phi^{-1}(1-\varepsilon))^2\, v + 2\tau\mu \leq \tau^2$.
        The value $1$ is excluded: for $c_1 = 0$ the tangent at $\tau$ is
        $\chi \leq 0$. At a Bernstein cell the tangent point must also satisfy
        $\mu_r < \tau - c_{i\ell}$ (Section~\ref{sec:consistency_checks}); otherwise
        no point with $v \geq 0$ satisfies the row (Section~\ref{sec:dyn_rainflow}).
  \item \texttt{pwl}: a piecewise tangent with $P = \texttt{pwl\_points}$ segments
        (default $8$). The breakpoints split $[0, \tau]$ into equal segments
        $[(p-1)\tau/P,\; p\tau/P]$ with midpoints $\hat\mu_p = (p - \tfrac12)\tau/P$,
        $p = 1, \ldots, P$. One binary $w_{i\ell kp}$ per cell, day and segment
        selects the segment that contains the mean:
        \begin{align}
          \textstyle\sum_{p=1}^{P} w_{i\ell kp} &= 1, \label{eq:a_pwl_sum} \\
          w_{i\ell kp} = 1 &\;\Rightarrow\;
            (p-1)\,\tau/P \leq \mu_{i\ell k} \leq p\,\tau/P, \label{eq:a_pwl_seg} \\
          w_{i\ell kp} = 1 &\;\Rightarrow\;
            \chi_{i\ell k} \leq f(\hat\mu_p) + f'(\hat\mu_p)\,(\mu_{i\ell k} - \hat\mu_p).
            \label{eq:a_pwl_cap}
        \end{align}
        The implications are Gurobi indicator constraints, whatever the value of
        \texttt{encoding}. With $P = 1$ the form is the single tangent at $\tau/2$,
        with no binaries.
\end{itemize}
At a cell whose reliability row is linear (Markov, Chernoff, inverse Gaussian, Gamma)
the row is built in its exact form, and the value of \texttt{reliability\_impl} is
ignored. An unknown value of \texttt{reliability\_impl} is an input error at every
cell (Section~\ref{sec:consistency_checks}).

\begin{proposition}[Safety of the linear forms]
\label{prop:a_tangent}
For every $\mu_r$ and every $\mu$,
$f(\mu_r) + f'(\mu_r)\,(\mu - \mu_r) = f(\mu) - c_2\,(\mu - \mu_r)^2 \leq f(\mu)$.
Therefore every point that satisfies the \texttt{tangent} or the \texttt{pwl} form
satisfies the \texttt{exact} form, and both linear forms lie inside the certificate
set $\mathcal{C}_{i\ell}$.
\end{proposition}
\begin{proof}
$f$ is a polynomial of degree $2$ in $\mu$ with leading coefficient $c_2$, so its
Taylor expansion at $\mu_r$ is exact: $f(\mu) = f(\mu_r) + f'(\mu_r)(\mu - \mu_r) +
c_2(\mu - \mu_r)^2$. Since $c_2 > 0$, the tangent lies below $f$. Each linear form
caps $\chi$ by one such tangent and keeps the side row. Therefore it implies
Eq.~\eqref{eq:a_quad}.
\end{proof}

\emph{Remark (accuracy).} The \texttt{tangent} form loses $c_2(\mu - \mu_r)^2$ in
$\chi$: it is exact at $\mu = \mu_r$ and loosest far from it. A tangent point near the
operating mean of the cell is the best choice. In the \texttt{pwl} form the loss is
$c_2(\mu - \hat\mu_p)^2 \leq c_2\,\tau^2/(4P^2)$. The midpoint of the segment that
contains $\mu$ is the nearest tangent point to $\mu$. So the \texttt{pwl} form
accepts exactly the points with $\chi \leq \max_p\, [f(\hat\mu_p) + f'(\hat\mu_p)(\mu -
\hat\mu_p)]$, the largest set that the $P$ tangents allow. The price is $P$ binaries
per cell and day.

\emph{Example (Cantelli, $\varepsilon = 0.1$, $\kappa_C^2 = 9$, $\tau = 1$, $v =
0.004$).} The largest admissible mean is $1 - \sqrt{9 \cdot 0.004} \approx 0.810$
with \texttt{exact}; $0.482$ with \texttt{tangent} at $0$ ($9v + 2\mu \leq 1$);
$0.714$ with \texttt{tangent} at $0.5\,\tau$ ($9v \leq 0.75 - \mu$); and $0.810$
with \texttt{pwl} and $P = 8$. The tangent at $0$ never admits a mean above
$\tau/2$, even with zero variance.

% ==== BLOCK 2 ====

% ----------------------------------------------------------------
\subsection{Repair Cost Variable}
\label{sec:repair_cost}
% ----------------------------------------------------------------

The variable $z_{i\ell k} \geq 0$ is the degradation removed by a repair of cell
$(i,\ell)$ on day $k$, in mean units, for every model and both laws. The objective
charges $C_{R,i\ell}\, z_{i\ell k}$ (Section~\ref{sec:objective}). The variable and
its rows are built only at cells with $C_{R,i\ell} > 0$, and only on the days whose
cost enters the objective: every day $k = 1, \ldots, T$, except the days of the
transitory phase in the mode \texttt{operating\_average}. The rows are
\begin{equation}
  z_{i\ell k} \;\geq\; \rho_{i\ell}\,\bigl(\mu_{i\ell,k-1} - g^{\mu}_{i\ell,k-1}\bigr)
    - \rho_{i\ell}\,\tau_{i\ell}\,\bigl(1 - x^m_{i\ell k}\bigr),
  \qquad z_{i\ell k} \geq 0,
  \label{eq:a_z}
\end{equation}
with $g^{\mu} \equiv 0$ under $\ARDinf$. On a repair day the first row reads $z_{i\ell
k} \geq \rho\,(\mu_{i\ell,k-1} - g^{\mu}_{i\ell,k-1})$, the removed amount of
Eq.~\eqref{eq:a_psi}. On any other day its right-hand side is at most $0$, because
$0 \leq \mu_{i\ell,k-1} - g^{\mu}_{i\ell,k-1} \leq \mu_{i\ell,k-1} \leq \tau$.

\emph{Replacement days.} A replacement day has $x^m_{i\ell k} = 0$, so $z_{i\ell k} =
0$ at the optimum. A replacement costs $C_{\mathrm{rep},i\ell}$ only. The replaced
part is discarded, so its damage does not change the price of the new part.

\emph{Remark (no upper rows).} The variable $z_{i\ell k}$ appears in no row other
than Eq.~\eqref{eq:a_z}, and in the objective with a positive coefficient. The
minimisation therefore drives it to the largest of its lower bounds:
$z_{i\ell k} = \rho\,(\mu_{i\ell,k-1} - g^{\mu}_{i\ell,k-1})$ on a repair day and
$z_{i\ell k} = 0$ on any other day. The upper rows of a McCormick envelope,
$z \leq \rho\,(\mu - g^{\mu})$ and $z \leq \rho\,\tau\,x^m$, would only remove points
that are never optimal, so they are not built. A solution that is not optimal (for
example, the best one found at the time limit) may carry a larger $z$; its reported
cost is then larger, never smaller.

\begin{proposition}[What $z$ measures]
\label{prop:a_repair_cost}
In every feasible solution, on every repair day $k$ of a cell, under both laws,
\[
  z_{i\ell k} \;\geq\; \rho\,\bigl(\mu_{k-1} - g^{\mu}_{k-1}\bigr)
  \;\geq\; \rho\,\bigl(\zeta^{\mu}_{k-1} - \zeta^{g}_{k-1}\bigr),
\]
where $\zeta^{\mu}$ and $\zeta^{g}$ are the mean and its anchor on the reference
trajectory (Section~\ref{sec:encoding}). When the starting anchor equals the true
mean level after the last intervention before the horizon, the right-hand side is
the expected degradation that the repair removes.
\end{proposition}
\begin{proof}
The first inequality is Eq.~\eqref{eq:a_z} with $x^m_{i\ell k} = 1$. Under $\ARDinf$,
$g^{\mu} = \zeta^{g} = 0$, and $\mu_{k-1} \geq \zeta^{\mu}_{k-1}$ by
Proposition~\ref{prop:a_domination}. Under $\ARDone$, let $t < k$ be the last
intervention day before $k$, or $t = 0$ if there is none. Then $g^{\mu}_{k-1} =
g^{\mu}_t$ and $\zeta^{g}_{k-1} = \zeta^{g}_t$.
\begin{itemize}[nosep,leftmargin=1.5em]
  \item If $t = 0$: $g^{\mu}_0 = \zeta^{g}_0$ (the starting anchor), so
        $g^{\mu}_{k-1} - \zeta^{g}_{k-1} = 0 \leq \mu_{k-1} - \zeta^{\mu}_{k-1}$.
  \item If $t \geq 1$: $g^{\mu}_t = \mu_t$ and $\zeta^{g}_t = \zeta^{\mu}_t$
        (Eq.~\eqref{eq:a_anchor}). On the days $j = t+1, \ldots, k-1$ there is no
        intervention: $\mu_j \geq \mu_{j-1} + \delta\mu_j$ and $\zeta^{\mu}_j =
        \zeta^{\mu}_{j-1} + \delta\mu_j$. So $\mu_j - \zeta^{\mu}_j$ does not
        decrease, and $\mu_{k-1} - \zeta^{\mu}_{k-1} \geq \mu_t - \zeta^{\mu}_t =
        g^{\mu}_{k-1} - \zeta^{g}_{k-1}$.
\end{itemize}
Rearranging gives $\mu_{k-1} - g^{\mu}_{k-1} \geq \zeta^{\mu}_{k-1} -
\zeta^{g}_{k-1}$. The mean recursion is linear (Section~\ref{sec:repair_laws}); with
an exact starting anchor, $\zeta^{\mu}$ and $\zeta^{g}$ are the true means, so
$\rho\,(\zeta^{\mu}_{k-1} - \zeta^{g}_{k-1}) = \mathbb{E}[D^- - D^+]$.
\end{proof}

A starting anchor above the true one, which the default may be
(Section~\ref{sec:repair_laws}), models a weaker repair. For the first repair before
any intervention of the horizon, the charged removal then follows the starting
anchor and may be below the true removal. From the first intervention on, the
reference gap $\zeta^{\mu} - \zeta^{g}$ equals the true gap, the sum of the mean
increments since that intervention, so the bound holds again.

Under both laws, then, the reported repair cost $\sum C_R\, z$ is an upper bound on
the expected repair cost of the schedule, with the exception above for a starting
anchor above the true one. Under $\ARDone$ the tracked mean and the
tracked anchor are both upper bounds, and their difference is still an upper bound
on the reference difference: the anchor carries the slack that the mean had at the
last intervention, and the slack of the mean does not decrease until the next one.
The bound is exact when the tracked mean has no slack.

% ----------------------------------------------------------------
\subsection{Damage Regularisation}
\label{sec:damage_reg}
% ----------------------------------------------------------------

The scalar $u \geq 0$ measures the largest damage of one train on one day, over the
whole horizon. The objective charges $C_D\,u$ (Section~\ref{sec:objective}). The
variable $u$ and its rows are built only when $C_D > 0$. With $C_D = 0$ they are not
built, and the input \texttt{penalty\_type} is ignored. The input check requires
$C_D \geq 0$ (Section~\ref{sec:consistency_checks}): a negative $C_D$ would reward
damage.

The input \texttt{penalty\_type} selects one of two forms. Both have one row per train
and day.
\begin{itemize}[leftmargin=1.5em]
  \item \texttt{inf\_norm} (default), a linear row:
        \begin{equation}
          u \;\geq\; \sum_{\ell=1}^{L} \mu_{i\ell k} \qquad \forall\, i, k .
          \label{eq:a_u_inf}
        \end{equation}
        At the optimum, $u = \max_{i,k} \sum_{\ell} \mu_{i\ell k}$.
  \item \texttt{quadratic}, a convex quadratic row:
        \begin{equation}
          u \;\geq\; \sum_{\ell=1}^{L} \mu_{i\ell k}^2 \qquad \forall\, i, k .
          \label{eq:a_u_quad}
        \end{equation}
        The row is a rotated second-order cone, $u \cdot 1 \geq \|\mu_{i\cdot k}\|_2^2$,
        which Gurobi solves without the parameter \texttt{NonConvex}. It requires
        \texttt{reliability\_impl = "exact"} at every cell
        (Section~\ref{sec:consistency_checks}): the linear forms are chosen for a
        model without quadratic rows, and this row is quadratic.
\end{itemize}

\emph{Role.} The term prices the worst train-level damage, so the minimisation
spreads damage across the fleet. It also removes slack from the tracked means where
$u$ is binding (Section~\ref{sec:encoding}). Since $\mu_{i\ell k} \geq
\zeta^{\mu}_{i\ell k} \geq 0$ (Proposition~\ref{prop:a_domination}), both forms of $u$
are upper bounds on the same quantity computed from the reference values.

% ----------------------------------------------------------------
\subsection{Loop Constraint}
\label{sec:loop}
% ----------------------------------------------------------------

The schedule of the operating phase is meant to repeat without end. One hard row per
tracked quantity and per anchor makes every repetition safe:
\begin{equation}
  s_{i\ell T} \;\leq\; s_{i\ell H_1}
  \qquad \text{for every cell } (i,\ell) \text{ and every coordinate } s
  \text{ of its full tracked state.}
  \label{eq:a_loop}
\end{equation}
The rows cover every tracked quantity of the cell (the mean, and the variance, $R$
or $K$ where the model tracks them) and, under $\ARDone$, every anchor. There is no
soft loop penalty in the objective.

\textbf{Endless schedule.} Repeat the operating phase without end: day $T + n H_2 +
h$, for $n \geq 0$ and $1 \leq h \leq H_2$, carries the binaries of day $H_1 + h$.
It also carries the increments of day $H_1 + h$, because the operating profile is
read cyclically with period $H_2$ (Section~\ref{sec:input_increments}). So every
repetition applies the same day maps $\varphi_{H_1+1}, \ldots, \varphi_T$. Write
$\Phi_h = \varphi_{H_1+h} \circ \cdots \circ \varphi_{H_1+1}$ for the first $h$
days of the operating phase and $\Phi = \Phi_{H_2}$ for the whole phase.

\begin{theorem}[Perpetual safety]
\label{thm:a_loop}
Let a solution satisfy every row of the problem. Then, for every cell and every day
$k \geq 1$ of the endless schedule, the true statistics $\theta_k$ and the reference
values $\zeta_k$ are at most a tracked value of the horizon that lies in
$\mathcal{C}_{i\ell}$. Therefore $\zeta_k \in \mathcal{C}_{i\ell}$ and
$\Pr(D_{i\ell k} > \tau_{i\ell}) \leq \varepsilon_{i\ell}$ on every day $k \geq 1$.
\end{theorem}
\begin{proof}
Fix a cell. Let $\xi_k$ be its tracked full state, and $\theta_k$ and $\zeta_k$ the
true statistics and the reference trajectory of the endless schedule. By (M2) and
Lemma~\ref{lem:a_monotone}, the induction of Proposition~\ref{prop:a_domination}
gives $\theta_k \leq \zeta_k$ on every day $k \geq 1$.
\begin{itemize}[nosep,leftmargin=1.5em]
  \item Days $k \leq T$: $\zeta_k \leq \xi_k$ (Proposition~\ref{prop:a_domination})
        and $\xi_k \in \mathcal{C}_{i\ell}$ (consequence on the horizon,
        Section~\ref{sec:encoding}).
  \item By Proposition~\ref{prop:a_domination} and Lemma~\ref{lem:a_monotone},
        $\xi_{H_1+h} \geq \varphi_{H_1+h}(\xi_{H_1+h-1}) \geq \cdots \geq
        \Phi_h(\xi_{H_1})$ for $h = 1, \ldots, H_2$. With Eq.~\eqref{eq:a_loop},
        $\Phi(\xi_{H_1}) \leq \xi_T \leq \xi_{H_1}$.
  \item Let $\Phi^n(\xi_{H_1}) \leq \xi_{H_1}$ (true for $n = 0$). Then
        $\Phi^{n+1}(\xi_{H_1}) = \Phi(\Phi^n(\xi_{H_1})) \leq \Phi(\xi_{H_1}) \leq
        \xi_{H_1}$ by Lemma~\ref{lem:a_monotone}. So $\Phi^n(\xi_{H_1}) \leq
        \xi_{H_1}$ for every $n \geq 0$.
  \item Day $k = H_1 + n H_2 + h$ with $n \geq 0$, $1 \leq h \leq H_2$:
        \[
          \zeta_k = \Phi_h\bigl(\Phi^n(\zeta_{H_1})\bigr)
          \leq \Phi_h\bigl(\Phi^n(\xi_{H_1})\bigr)
          \leq \Phi_h(\xi_{H_1})
          \leq \xi_{H_1+h} \in \mathcal{C}_{i\ell}.
        \]
\end{itemize}
Therefore, on every day $k \geq 1$, $\theta_k \leq \zeta_k \leq s$ for a tracked value
$s \in \mathcal{C}_{i\ell}$. Since $\mathcal{C}_{i\ell}$ is a down-set,
$\zeta_k \in \mathcal{C}_{i\ell}$, and (M3) gives $\Pr(D_{i\ell k} > \tau_{i\ell})
\leq \varepsilon_{i\ell}$.
\end{proof}

The theorem holds for every model and both repair laws, because it uses only
the facts (M1)--(M3) and Lemma~\ref{lem:a_monotone}. It holds in both encodings and
in every reliability form. It needs no relation between $H_1$ and $H_2$: the
transitory phase may be shorter or longer than the operating phase, because
$H_1$ only marks the day at which the repeated block starts. Slack in the tracked
values does not weaken it: the proof compares the reference trajectory with the
tracked values, whatever slack they carry.

\emph{Remark (why the loop row covers the full state).} The third step of the proof
needs $\Phi(\xi_{H_1}) \leq \xi_{H_1}$ in every coordinate, anchors included.
Example: law $\ARDone$, $\rho = 0.5$, the first phase ends at $(\mu, g) = (0.4,
0.1)$, and the operating phase is a mission ($+0.1$), a repair, a mission
($+0.05$). The phase ends at $(0.35, 0.3)$, since the repair gives $0.5 \cdot 0.5 +
0.5 \cdot 0.1 = 0.3$. The row on the mean holds ($0.35 \leq 0.4$); the row on the
anchor fails ($0.3 > 0.1$). Each repetition raises the state after the repair by
$0.075$, so the mean grows without bound: the third repetition already peaks at
$0.525$, above the checked peak $0.5$. A loop row on the mean alone would accept
this unsafe schedule.

\emph{Remark (no loop row on a scalar bound).} The loop row must compare the tracked
quantities one by one. A single comparison of a scalar function of them, such as the
failure-probability bound $U(\mu, v)$, does not chain: $U$ takes the same value at
many states, and its level sets are not invariant under the day maps. Example
(Cantelli, $U = v / (v + (\tau - \mu)^2)$, $\tau = 10$, $\rho = 0.5$, law
$\ARDinf$; the operating phase is one repair day followed by mission days that add
$3$ to the mean and $3.99$ to the variance in total): from $(\mu, v) = (4, 9)$,
$U = 0.2000$ at the end of the first phase and $U = 0.1997$ at the end of the
operating phase, so a scalar row would accept. The next repetition gives $U =
0.2151 > 0.2$, and $U$ tends to $0.2495$. Eq.~\eqref{eq:a_loop} rejects the
schedule, since the mean rises from $4$ to $5$. The loop row states an order between
tracked quantities; it is not a stochastic order between damage laws.

\begin{proposition}[Replacement under $\ARDone$]
\label{prop:a_ard1_replacement}
Let a cell follow $\ARDone$ with $\rho_{i\ell} < 1$, and let its train fly, in the
operating phase, at least one mission with a positive mean increment for the cell.
Then every feasible solution has at least one replacement of the cell in the
operating phase.
\end{proposition}
\begin{proof}
Let the cell have no replacement in the operating phase. By the anchor rows,
$g_k = g_{k-1}$ on a day without intervention and $g_k = \mu_k$ on a repair day.
Every day satisfies $\mu_k \geq g_k$: on an intervention day with equality, and on
other days because $\mu$ does not decrease and $g$ stays constant (with $g_0 \leq
\mu_0$, Section~\ref{sec:repair_laws}). On a repair day $k$ of the operating phase,
\[
  g_k = \mu_k \geq (1-\rho)\,\mu_{k-1} + \rho\, g_{k-1}
  = g_{k-1} + (1-\rho)\,(\mu_{k-1} - g_{k-1}).
\]
So $g_T \geq g_{H_1} + (1-\rho) \sum (\mu_{k-1} - g_{k-1})$, the sum over the
repair days of the operating phase, each term non-negative. The anchor row $g_T
\leq g_{H_1}$ and $\rho < 1$ force every term to $0$. Then every repair day gives
$\mu_k \geq \mu_{k-1}$, every other day gives $\mu_k \geq \mu_{k-1} + \delta\mu_k$,
and
\[
  \mu_T \geq \mu_{H_1} + \textstyle\sum_{k = H_1+1}^{T} \delta\mu_k > \mu_{H_1},
\]
which violates the row $\mu_T \leq \mu_{H_1}$. Therefore the cell is replaced at
least once in the operating phase.
\end{proof}

The requirement is not an artefact of the rows. Without a replacement, each repair
leaves the share $1-\rho$ of the damage since the last intervention in place, so the
reference mean grows without bound over the repetitions, as in the first remark
above. With $\rho = 1$ a repair returns the state to the anchor, and no replacement
is needed. The Wiener model needs a replacement in the operating phase for another
reason, its variance (Section~\ref{sec:dyn_wiener}).
